\documentclass[10pt,a4paper]{amsart}
\usepackage[margin=19mm]{geometry}
\usepackage{amsmath,amssymb,mathtools}
\usepackage[scr=rsfs,cal=euler]{mathalfa}
\usepackage{xfrac}
\usepackage[unicode,hidelinks,bookmarks=false]{hyperref}
\newcommand{\ie}{i.e.\ }
\newcommand{\cf}{cf.\ }
\newcommand{\resp}{respectively\ }
\newcommand{\Subsection}{\subsection}
\providecommand{\nobreakdash}{\nobreak\hbox{-}}
\setlength{\emergencystretch}{2em}
\multlinegap0pt
\usepackage{amssymb}
\usepackage{float}
\usepackage[all]{xy}
\usepackage{tikz}
\newtheorem{proposition}{Proposition}
\newtheorem{lemma}{Lemma}
\newcommand{\Whitney}{\rmW}
\newcommand{\rmL}{\mathrm{L}}
\newcommand{\rmW}{\mathrm{W}}
\newcommand{\rmd}{\mathrm{d}}
\title{Interpretation of a Discrete de Rham method \\ as a Finite Element System}
\author{Snorre H. Christiansen, Francesca Rapetti}
\date{Source: arXiv:2512.05912v2 (CC BY 4.0)}
\begin{document}
\maketitle

\noindent\emph{Source and licence.} Interpretation of a Discrete de Rham method \\ as a Finite Element System. Version arXiv:2512.05912v2 (CC BY 4.0), distributed by arXiv under the Creative Commons Attribution 4.0 International licence, http://creativecommons.org/licenses/by/4.0/, as recorded in the arXiv licence metadata for this exact version and shown on the abstract page https://arxiv.org/abs/2512.05912v2. Full licence notice: \texttt{licenses/arxiv-2512.05912-CC-BY-4.0.txt}.\par\smallskip

\noindent\textit{Editorial scope.} Counted unit: the article's proposition proposition (source0.tex lines 1124--1125). The article's complete original proof of it is delivered with it (source0.tex lines 1127--1160, one \texttt{proof} environment of 34 lines). No mathematics is written, repaired or re-derived: the statement and the proof are the article's own lines, in the article's own notation, and no retained line is substituted or re-flowed.  Retained uncounted as the source-local context the proof needs: lines 1174--1178; lines 1193--1198.  Nothing was omitted from any retained span, so the package carries no omission record.  No retained line contains a citation, so the wrapper carries no bibliography and no bibliography entry.  No retained line contains a figure environment, an image-inclusion command or drawing code, so no image is delivered and the package declares no asset dependency.  Editorial formatting only, in addition: an AMS \texttt{amsart} preamble that restates the article's own declarations and macros used by the retained lines, namely the environments \texttt{proposition}, \texttt{lemma} on separate counters, together with the standard packages those lines need.  The retained environments print in this document as Proposition 1, Lemma 1 in order of appearance, whereas the article prints the same items with the numbering of its own full document.  The complete native source of the article as distributed in the arXiv e-print for this exact version is bundled unchanged as \texttt{sources/190447/source0.tex}.
\begin{lemma}[Inf-Sup condition]\label{lem:infsup} Fix $k$ and $T$. Suppose that $E$ is a finite dimensional subspace of $X^k(T)$  and that $E \cap \rmd Y^{k-1}(T) = 0$. Then:
\begin{equation}
  \inf_{v \in \rmd Y^{k-1}(T), w \in E} \sup_{u\in X^k_0(T)} \frac{|\langle u, v \rangle + \langle u, w \rangle|}{\| u\| (\| v \| + \| w\|)} >0.
\end{equation}
\end{lemma}

\begin{lemma}[Coercivity on the kernel]\label{lem:coerker}
  There is $C>0$ such that for all $u \in Z^k_0(T)$, such that $u \perp \rmd Y^{k-1}(T)$ (w.r.t. $\rmL^2$ scalar product):
\begin{equation}
   |u| \leq C | \rmd u|.
\end{equation}
\end{lemma}

\begin{proposition} On a cell $T$ of dimension $d$ and for $k < d$, the above system of PDEs for $B^k(T)$ is wellposed, for given boundary data.
\end{proposition}

\begin{proof}
Extend the boundary data to an arbitrary field $\tilde u \in X^k(T)$ (for instance by ordinary harmonic extension). Then write the unknown $u \in X^k(T)$ as $u = \tilde u + \hat u$ with $\hat u \in X^k_0(T)$. We now consider the system:
\begin{align}
  \int_T \rmd \hat u \cdot \rmd u' + \int_T \rmd u' \wedge \kappa_T a + \int_T u' \wedge \kappa_T b + \int_T u' \cdot \rmd v & = - \int_T \rmd \tilde u \cdot \rmd u',\\
  \int_T \rmd \hat u \wedge \kappa_T a' = - \int_T \rmd \tilde u \wedge \kappa_T a'  ,\ \int_T \hat u \wedge \kappa_T b' & = - \int_T \tilde u \wedge \kappa_T b' ,\\
  \int_T \hat u \cdot \rmd v' & =  - \int_T \tilde u \cdot \rmd v' ,
\end{align}


Since $k < d$ we can now integrate by parts the constraints, as in:
\begin{align}
\int_T \rmd \hat u \wedge \kappa_T a' + \int_T \hat u \wedge \kappa_T b' & = \int_T \hat u \wedge ( (-1)^{k+1} \rmd \kappa_T a' + \kappa_T b'),\\
&=  \int_T \hat u \wedge (\frac{(-1)^{k+1}}{d-k} a' + \kappa_T b') ,
\end{align}
For $k = 0$ there is no $b'$ term. %For $k = d$ there is no $a'$ term.
We recognize the general expression of a $w' \in \Whitney^{d-k}(T)$:
\begin{equation}\label{eq:wprime}
  w' = \frac{(-1)^{k+1}}{d-k} a' + \kappa_T b',
\end{equation}
uniquely corresponding to $a'$ and $b'$.

We can rewrite the system for  $u \in X^k_0(T)$ as, for all $u' \in X^k_0(T)$:
\begin{align}
  \int_T \rmd \hat u \cdot \rmd u' + \int_T u' \wedge w + \int_T u' \cdot \rmd v & = - \int_T \rmd \tilde u \cdot \rmd u',\\
  \int_T  \hat u \wedge w' & = - \int_T \rmd \tilde u \wedge \kappa_T a' - \int_T \tilde u \wedge \kappa_T b', \\
  \int_T \hat u \cdot \rmd v' & =  - \int_T \tilde u \cdot \rmd v' ,
\end{align}
for some $w \in W^{d-k}(T)$ and $v \in Y^{k-1}(T)$, and for all $w' \in W^{d-k}(T)$ (written as (\ref{eq:wprime}) ) and $v' \in Y^{k-1}(T)$.


We look first at the homogeneous problem (zero right sides). The second equation says that $\hat u \in Z^k_0(T)$ and the third one then says that $\hat u\in Y^k(T)$. The first equation with $u' = \hat u$ gives $\rmd \hat u =0$. This shows that $\hat u = 0$. The homogeneous problem thus has only the trivial solution.

The proof of the well-posedness in the PDE sense, is contained in the Lemmas \ref{lem:infsup} and \ref{lem:coerker}, below.
\end{proof}

\end{document}
